\subsection{適用条件}

極限を計算していると,分子も分母も$0$へ近づく商が現れる.高校の範囲では,このような極限を微分係数の定義に帰着させて求めることが多い.実際,\,$f(c)=g(c)=0$で,\,$f,g$が$c$で微分可能,\,$g'(c)\ne0$なら,分子と分母を$x-c$で割ってやれば
\[
 \frac{f(x)}{g(x)}
 =\cfrac{\cfrac{f(x)-f(c)}{x-c}}{\cfrac{g(x)-g(c)}{x-c}}
 \longrightarrow\frac{f'(c)}{g'(c)}
 \qquad(x\to c)
\]
となる.例えば$\lim_{x\to0}\frac{\e^x-1}{\sin x}=\frac11=1$である.

しかし,分母が$x^3$のように$g'(c)=0$となる場合には,この方法は使えない.また,\,$x\to\infty$のように$c$が有限でない場合や,分子と分母がともに無限大へ発散する場合も扱えない.そこで,微分係数の値ではなく,分子と分母の導関数の比の極限を調べることを考える.ただし,元の商の形を確かめず,機械的に分子と分母を微分してはいけない.次の条件を確かめる必要がある.

\begin{theorem}{ロピタルの定理の一側極限版}
$x$が区間の一端$c$へ近づく一側極限を考える.\,$c$は有限でも$+\infty,-\infty$でもよい.端点を除くその区間で$f,g$は微分可能であり,\,$g'(x)\ne0$とする.さらに,次のいずれかを満たすとする.
\begin{enumerate}[label=(\roman*),leftmargin=*]
 \item $f(x)\to0$かつ$g(x)\to0$である.
 \item $|f(x)|\to\infty$かつ$|g(x)|\to\infty$である.
\end{enumerate}
このとき
\[
 \lim_{x\to c}\frac{f'(x)}{g'(x)}=L
\]
が存在すれば
\[
 \lim_{x\to c}\frac{f(x)}{g(x)}=L
\]
である.ここで$L$は実数のほか,\,$+\infty,-\infty$でもよい.元の商は端点の十分近くで定義されている.
\end{theorem}

両側極限を求めるときは,左右の一側極限にそれぞれ適用し,結果が一致することを確かめる.また,\,$\frac{f'}{g'}$にもう一度定理を使いたい場合は,その新しい商についても同じ条件を改めて確かめる.微分後の商が$0$と$0$の商になるとは限らない.
